\section{Critical-height relationships on the same rectangles}
\label{f:sec:compatible-rectangles}

A finite rectangle is the velocity--pressure--force record selected at fixed
temporal and spatial orders from the one forced history.  Ranging over the
mixed indices exhausts all of these finite records.  Independently,
differentiating the critical heights and squaring the viscous cell produces relationships on those
same rectangle coordinates.  Their completed heights recover every spatial
Sobolev row, and the forced recurrence reconstructs the same temporal and
pressure intersection, with the nonlinear, pressure, and force faces retained.


\begin{proposition}[Kinetic-energy identity]
\label{f:prop:kinetic-energy}
For every \(0\le t<T_*\),
\begin{equation}
 \frac12\norm{u(t)}_2^2
 +\nu\int_0^t\norm{\nabla u(s)}_2^2\dd s
 =\frac12\norm{u_0}_2^2
 +\int_0^t\ip{f(s)}{u(s)}_{L^2}\dd s.
 \label{f:eq:kinetic-energy}
\end{equation}
Put
\[
 K_f(T):=\norm{u_0}_2+\int_0^T\norm{\Pp f(s)}_2\dd s.
\]
Consequently, for \(0<T<T_*\),
\begin{align}
 \norm u_{L^\infty(0,T;L^2)}&\le K_f(T),
 \label{f:eq:energy-Linf}\\
 \norm u_{L^2(0,T;H^1)}^2
 &\le\left(T+\frac1{2\nu}\right)K_f(T)^2.
 \label{f:eq:energy-L2H1}
\end{align}
\end{proposition}

\begin{proof}
Take the \(L^2\) pairing of the momentum equation with \(u\).  Integration
by parts and \(\nabla\cdot u=0\) give
\[
 \ip{(u\cdot\nabla)u}{u}_{L^2}=0,
 \qquad
 \ip{\nabla p}{u}_{L^2}=0,
 \qquad
 \ip{\Delta u}{u}_{L^2}=-\norm{\nabla u}_2^2.
\]
Since \(u\) is divergence free,
\(\ip f u=\ip{\Pp f}u\).  Integration in time proves
\eqref{f:eq:kinetic-energy}.  Applying the identity first to
\((\norm u_2^2+\varepsilon)^{1/2}\) and then letting
\(\varepsilon\downarrow0\) gives \(\norm{u(t)}_2\le K_f(t)\).  Moreover,
\[
 \int_0^T K_f(t)\norm{\Pp f(t)}_2\dd t
 =\frac12\bigl(K_f(T)^2-\norm{u_0}_2^2\bigr),
\]
so \eqref{f:eq:kinetic-energy} gives
\(\int_0^T\norm{\nabla u}_2^2\dd t\le K_f(T)^2/(2\nu)\).
Adding \(\int_0^T\norm{u(s)}_2^2\dd s\) proves
\eqref{f:eq:energy-L2H1}.
\end{proof}

Pairing a differentiated velocity equation with a spatial derivative of
that row expresses the change of its Sobolev energy through viscous
dissipation, the full transport--pressure contribution, and the prescribed
force work. Repeated time differentiation then exposes higher spatial energies,
with every differentiated nonlinear and force term retained in each identity.

\begin{definition}[Completed height rows]
\label{f:def:critical-height-rows}
For \(r\in\Nzero\) and \(s\ge0\), set
\begin{align}
 E_{r,s}(t)&:=\frac12\norm{\Lambda^sV_r(t)}_2^2,
 \label{f:eq:temporal-row-spatial-energy}\\
 \cP_{r,s}(t)&:=-\operatorname{Re}
 \ip{\Lambda^sV_r(t)}
 {\Lambda^s\!\left(
   \sum_{a=0}^{r}\binom ra(V_a\cdot\nabla)V_{r-a}+\nabla Q_r
 \right)(t)}_{L^2}.
 \label{f:eq:temporal-row-pressure-complete-transfer}
 \\
 \mathcal W_{r,s}(t)&:=\operatorname{Re}
 \ip{\Lambda^sV_r(t)}{\Lambda^sF_r(t)}_{L^2}.
 \label{f:eq:temporal-row-force-work}
\end{align}
For the base row write \(E_s:=E_{0,s}\), \(\cP_s:=\cP_{0,s}\),
\(\mathcal W_s:=\mathcal W_{0,s}\), and \(H:=E_{1/2}\).
\end{definition}

Differentiating this quadratic height along the classical history gives
\[
 H^{(r)}(t)=\frac12\sum_{a=0}^r\binom ra
 \ip{\Lambda^{1/2}V_a(t)}{\Lambda^{1/2}V_{r-a}(t)}_{L^2}.
\]
On every strict interval the mixed jet makes this contraction an identity of
the same forced history.  The following recurrence exposes its spatial face
without discarding the transfer or force work.  For each finite \(k\), choose
a coordinate rectangle high enough to contain the required differentiated
rows.  Deriving the height identities on those rows gives a finite relationship
rectangle containing \(H^{(k)}\), every differentiated transfer and force-work
term, and the spatial energy \(E_{k+1/2}\).  Compatibility under enlargement
turns these finite identities into the completed-height exhaustion below.

\begin{proposition}[Completed height recurrence in every temporal row]
\label{f:prop:completed-height-recurrence}
For every strict classical interval \((0,T)\), every \(r\in\Nzero\), and
every \(s\ge0\),
\begin{equation}
 \partial_tE_{r,s}=\cP_{r,s}+\mathcal W_{r,s}-2\nu E_{r,s+1}.
 \label{f:eq:sobolev-height-balance}
\end{equation}
For every \(k\in\N\),
\begin{equation}
 \partial_t^kE_{r,1/2}
 =(-2\nu)^kE_{r,k+1/2}
 +\sum_{j=0}^{k-1}(-2\nu)^j
   \partial_t^{k-1-j}
   (\cP_{r,j+1/2}+\mathcal W_{r,j+1/2}).
 \label{f:eq:completed-height-every-temporal-row}
\end{equation}
In the base row \(r=0\),
\begin{equation}
 H^{(k)}
 =(-2\nu)^kE_{k+1/2}
 +\sum_{j=0}^{k-1}(-2\nu)^j
   \partial_t^{k-1-j}(\cP_{j+1/2}+\mathcal W_{j+1/2}),
 \label{f:eq:completed-height-recurrence}
\end{equation}
and the highest spatial energy is the exact coordinate
\begin{equation}
 E_{k+1/2}
 =\frac{H^{(k)}-
   \displaystyle\sum_{j=0}^{k-1}(-2\nu)^j
   \partial_t^{k-1-j}(\cP_{j+1/2}+\mathcal W_{j+1/2})}
 {(-2\nu)^k}.
 \label{f:eq:completed-height-row}
\end{equation}
\end{proposition}

\begin{proof}
Pair the temporal recurrence for \(V_r\) with \(\Lambda^{2s}V_r\).
The pressure-complete nonlinear sum is exactly \(\cP_{r,s}\), the force
pairing is exactly \(\mathcal W_{r,s}\), and
\[
 \nu\operatorname{Re}\ip{\Lambda^sV_r}{\Lambda^s\Delta V_r}_{L^2}
 =-\nu\norm{\Lambda^{s+1}V_r}_2^2=-2\nu E_{r,s+1}.
\]
This proves \eqref{f:eq:sobolev-height-balance}.  Induction on \(k\) gives
the stronger identity
\begin{equation}
 \partial_t^kE_{r,s}
 =(-2\nu)^kE_{r,s+k}
 +\sum_{j=0}^{k-1}(-2\nu)^j
   \partial_t^{k-1-j}(\cP_{r,s+j}+\mathcal W_{r,s+j}).
 \label{f:eq:all-height-recurrence}
\end{equation}
Indeed, differentiation of the order-\(k\) identity followed by
\(\partial_tE_{r,s+k}=\cP_{r,s+k}+\mathcal W_{r,s+k}
-2\nu E_{r,s+k+1}\) produces the
order-\(k+1\) identity.  Taking \(s=1/2\) proves
\eqref{f:eq:completed-height-every-temporal-row}.  Taking \(r=0\) gives
\eqref{f:eq:completed-height-recurrence}; since \(\nu>0\), division by
\((-2\nu)^k\) gives \eqref{f:eq:completed-height-row}.
\end{proof}

\begin{definition}[Completed critical-height coordinates]
\label{f:def:completed-critical-height-coordinates}
Set
\begin{equation}
 R_0:=H=E_{1/2},
 \qquad
 R_k
 :=\frac{H^{(k)}-
   \displaystyle\sum_{j=0}^{k-1}(-2\nu)^j
   \partial_t^{k-1-j}(\cP_{j+1/2}+\mathcal W_{j+1/2})}
 {(-2\nu)^k}
 \quad(k\in\N).
 \label{f:eq:completed-critical-coordinate}
\end{equation}
Proposition~\ref{f:prop:completed-height-recurrence} gives the exact identity
\begin{equation}
 R_k=E_{k+1/2}
 =\frac12\norm{\Lambda^{k+1/2}u}_2^2
 \qquad(k\in\Nzero).
 \label{f:eq:completed-critical-coordinate-identity}
\end{equation}
\par
\end{definition}

\begin{proposition}[Complete forced rectangle cell]
\label{f:prop:complete-forced-rectangle-cell}
For every temporal row \(n\in\Nzero\), every \(s\ge0\), and every strict
classical interval, the projected differentiated equation is
\begin{equation}
 \Lambda^sV_{n+1}+\nu\Lambda^{s+2}V_n
 =\Lambda^s\Pp(F_n-B_n),
 \qquad
 B_n:=\sum_{a=0}^n\binom na(V_a\cdot\nabla)V_{n-a}.
 \label{f:eq:complete-forced-rectangle-cell}
\end{equation}
Its square is the viscous diagonal identity
\begin{equation}
 \norm{\Lambda^sV_{n+1}}_2^2
 +\nu^2\norm{\Lambda^{s+2}V_n}_2^2
 +\nu\frac{\dd}{\dd t}\norm{\Lambda^{s+1}V_n}_2^2
 =\norm{\Lambda^s\Pp(F_n-B_n)}_2^2.
 \label{f:eq:complete-forced-rectangle-square}
\end{equation}
Before projection, the gradient face is orthogonal to both divergence-free
faces.  Consequently,
\begin{align}
 &\norm{V_{n+1}}_{\dot H^s}^2
 +\nu^2\norm{\Delta V_n}_{\dot H^s}^2
 +\norm{\nabla Q_n}_{\dot H^s}^2
 +\nu\frac{\dd}{\dd t}\norm{\nabla V_n}_{\dot H^s}^2
 \notag\\
 &\hspace{4cm}=\norm{F_n-B_n}_{\dot H^s}^2.
 \label{f:eq:complete-forced-pressure-square}
\end{align}
Thus, for \(0<S<T_*\),
\begin{align}
 &\nu\norm{\Lambda^{s+1}V_n(S)}_2^2
 +\int_0^S\!\left(
   \norm{\Lambda^sV_{n+1}}_2^2
   +\nu^2\norm{\Lambda^{s+2}V_n}_2^2
   +\norm{\Lambda^s\nabla Q_n}_2^2\right)\dd t
 \notag\\
 &\qquad=
 \nu\norm{\Lambda^{s+1}V_n^0}_2^2
 +\int_0^S\norm{\Lambda^s(F_n-B_n)}_2^2\dd t.
 \label{f:eq:complete-forced-rectangle-integrated}
\end{align}
These identities retain the temporal, viscous, pressure, transport, and force
faces of the differentiated momentum row.
\end{proposition}

\begin{proof}
Apply \(\Pp\Lambda^s\) to the temporal row in
Proposition~\ref{f:prop:formal-mixed-recurrences}. Since
\(-\Delta=\Lambda^2\), moving the viscous term to the left gives
\eqref{f:eq:complete-forced-rectangle-cell} with a positive sign.
On squaring, the cross term is
\[
 2\nu\operatorname{Re}
 \ip{\Lambda^sV_{n+1}}{\Lambda^{s+2}V_n}_{L^2}
 =\nu\frac{\dd}{\dd t}\norm{\Lambda^{s+1}V_n}_2^2,
\]
which proves \eqref{f:eq:complete-forced-rectangle-square}.  The complementary
projection gives
\(\nabla Q_n=(I-\Pp)(F_n-B_n)\).  Orthogonality of \(\Pp\) and
\(I-\Pp\) proves \eqref{f:eq:complete-forced-pressure-square}; integration in
time gives \eqref{f:eq:complete-forced-rectangle-integrated}.
\end{proof}

\Needspace{14\baselineskip}
\begin{proposition}[Complete-equation estimates on finite rectangles]
\label{f:prop:datum-complete-rectangle-estimates}
Assume the hypotheses of Theorem~\ref{f:thm:datum-rectangles}.  The rectangle
identities estimate the complete equation inside each selected finite rectangle.
For every \(n\in\Nzero\) and \(s\ge0\), put
\begin{align}
 \mathfrak D_{n,s}(S):={}&
 \nu\norm{\Lambda^{s+1}V_n(S)}_2^2
 +\int_0^S\!\left(
   \norm{\Lambda^sV_{n+1}}_2^2
   +\nu^2\norm{\Lambda^{s+2}V_n}_2^2
   +\norm{\Lambda^s\nabla Q_n}_2^2\right)\dd t.
 \label{f:eq:datum-complete-row-record}
\end{align}
Then
\begin{equation}
 \sup_{0<S<T_*}\mathfrak D_{n,s}(S)
 \le
 \nu\norm{\Lambda^{s+1}V_n^0}_2^2
 +\norm{\Lambda^s(F_n-B_n)}_{L^2(0,T_*;L^2)}^2
 <\infty.
 \label{f:eq:datum-complete-row-projection-bound}
\end{equation}
Thus the selected rectangle bounds, on the same terminal interval, its
adjacent temporal face, viscous face, full binomial nonlinear row, normalized
pressure-gradient face, and prescribed-force face.
\end{proposition}

\begin{proof}
Let \(I=(0,T_*)\).  Fix \((n,s)\).  Choose a datum-generated finite rectangle high enough to
contain every velocity factor in \(B_n\) at spatial order \(s+2\).
The adjacent trace lemma gives the required \(L^\infty_tH^{s+2}_x\) factor,
while the higher spatial face gives the other factor in
\(L^2_tH^{s+2}_x\).  Sobolev multiplication, summed over the finite binomial
row, gives the explicit same-history estimate
\begin{equation}
 \norm{B_n}_{L^2(I;\dot H^s)}
 \le C_s\sum_{a=0}^{n}\binom na
 \norm{V_a}_{L^\infty(I;H^{s+2})}
 \norm{V_{n-a}}_{L^2(I;H^{s+2})}
 <\infty.
 \label{f:eq:datum-complete-binomial-projection-bound}
\end{equation}
The prescribed Schwartz force gives
\(F_n\in L^2(I;\dot H^s)\).  Now evaluate the independently derived
positive-sign identity \eqref{f:eq:complete-forced-rectangle-integrated} on
these selected coordinates.  Its right-hand side is finite, and taking the
supremum over \(S<T_*\) proves
\eqref{f:eq:datum-complete-row-projection-bound}.  The pressure term is present
because \eqref{f:eq:complete-forced-pressure-square} is the unprojected
orthogonal completion of the same row.  Every quantity in the resulting
estimate is a coordinate of that forced equation row.
\end{proof}

\begin{proposition}[Complete momentum rows inside one finite rectangle]
\label{f:prop:rectangle-bounds-complete-rows}
Assume the hypotheses of Theorem~\ref{f:thm:datum-rectangles}.  Fix
\(N\in\Nzero\), \(M\ge2\), and \(0<T\le T_*\), and put
\begin{align*}
 R_v&:=\sum_{n=0}^N\left(
 \norm{V_n}_{L^2(0,T;H^{M+1})}^2
 +\norm{V_{n+1}}_{L^2(0,T;H^{M-1})}^2\right),\\
 I_{N,M}&:=\sum_{n=0}^N\norm{V_n^0}_{H^M}^2.
\end{align*}
Then the adjacent faces of this same rectangle give
\begin{equation}
 \sum_{n=0}^N\sup_{0\le t\le T}\norm{V_n(t)}_{H^M}^2
 \le I_{N,M}+R_v,
 \label{f:eq:rectangle-adjacent-trace-bound}
\end{equation}
and the complete binomial nonlinear rows obey
\begin{equation}
 \sum_{n=0}^N\norm{B_n}_{L^2(0,T;H^{M-1})}^2
 \le C_{N,M}(I_{N,M}+R_v)R_v.
 \label{f:eq:rectangle-complete-binomial-bound}
\end{equation}
For \(0<S\le T\), define the pressure-complete dissipation record
\begin{align*}
 \mathcal D_{N,M}(S):=\sum_{n=0}^N\bigg[&
 \nu\norm{V_n(S)}_{\dot H^M}^2
 +\norm{V_{n+1}}_{L^2(0,S;\dot H^{M-1})}^2\\
 &+\nu^2\norm{V_n}_{L^2(0,S;\dot H^{M+1})}^2
 +\norm{\nabla Q_n}_{L^2(0,S;\dot H^{M-1})}^2\bigg].
\end{align*}
The integrated complete rows satisfy
\begin{align}
 \mathcal D_{N,M}(S)
 \le{}&\nu I_{N,M}
 +2\sum_{n=0}^N\norm{F_n}_{L^2(0,T;H^{M-1})}^2
 \notag\\
 &+2C_{N,M}(I_{N,M}+R_v)R_v.
 \label{f:eq:rectangle-complete-row-bound}
\end{align}
Thus one finite velocity rectangle, together with the displayed prescribed
force coordinates, bounds on the same interval and in the same equations its
adjacent temporal rows, viscous Laplacians, complete nonlinear rows, and
normalized pressure gradients.  Lower spatial orders follow by Sobolev nesting
from a rectangle with \(M\ge2\).
\end{proposition}

\begin{proof}
For \(0\le n\le N\), the Hilbert triple
\(H^{M+1}\subset H^M\subset H^{M-1}\) and
\(\partial_tV_n=V_{n+1}\) give
\[
 \norm{V_n(t)}_{H^M}^2-\norm{V_n^0}_{H^M}^2
 =2\operatorname{Re}\int_0^t
 \ip{(1-\Delta)^{(M+1)/2}V_n}
 {(1-\Delta)^{(M-1)/2}V_{n+1}}_{L^2}\dd s.
\]
Cauchy--Schwarz in time and \(2ab\le a^2+b^2\), followed by summation in
\(n\), prove \eqref{f:eq:rectangle-adjacent-trace-bound}.  Sobolev
multiplication applied to every allocation in \(B_n\), using one factor in
\(L^\infty(0,T;H^M)\) and the other in
\(L^2(0,T;H^{M+1})\), proves
\eqref{f:eq:rectangle-complete-binomial-bound}.  Finally, sum
\eqref{f:eq:complete-forced-rectangle-integrated} over \(0\le n\le N\) at
\(s=M-1\), use
\(\norm{F_n-B_n}^2\le2\norm{F_n}^2+2\norm{B_n}^2\), and substitute
\eqref{f:eq:rectangle-complete-binomial-bound}.  This gives
\eqref{f:eq:rectangle-complete-row-bound}.  The identities
\[
 \nabla Q_n=(I-\Pp)(F_n-B_n),
 \qquad
 V_{n+1}=\nu\Delta V_n-\Pp B_n+\Pp F_n
\]
show directly that these are the terms of the same pressure-complete row.
\end{proof}

The finite sum above evaluates the coordinates selected by
Definition~\ref{f:def:finite-rectangle}.  Theorem~\ref{f:thm:datum-rectangles}
records their finite norm, while the complete-equation estimate retains and
relates every face of those coordinates.  The next proposition records both
directions between adjacent-row membership and its finite rectangle norm,
together with monotone exhaustion of the interval.

\Needspace{14\baselineskip}
\begin{proposition}[Norm equivalence and full-interval realization]
\label{f:prop:datum-terminal-realization}

For the fixed maximal history, the family of bounds
\eqref{f:eq:datum-rectangle-bound} is equivalent to the adjacent-row statement
\eqref{f:eq:datum-adjacent-rows}.  Whenever either formulation holds, for every
\(N,M\in\Nzero\),
\begin{equation}
 \sum_{n=0}^{N}\left(
 \norm{V_n}_{L^2(0,T_*;H^{M+1})}^2
 +\norm{V_{n+1}}_{L^2(0,T_*;H^{M-1})}^2
 \right)
 =\lim_{T\uparrow T_*}\mathfrak R_{N,M}(u;T)
 =\mathfrak R_{N,M}^*[u_0,\nu,f,T_*].
 \label{f:eq:datum-terminal-monotone-realization}
\end{equation}
Every row is the corresponding distributional derivative of the single
pressure-normalized maximal history.
\end{proposition}

\begin{proof}
Assume first \eqref{f:eq:datum-adjacent-rows}.  For fixed \(N,M\), set
\[
 C_{N,M}:=
 \sum_{n=0}^{N}\left(
  \norm{V_n}_{L^2(0,T_*;H^{M+1})}^2
  +\norm{V_{n+1}}_{L^2(0,T_*;H^{M-1})}^2
 \right).
\]
The sum is finite because it contains finitely many instances of
\eqref{f:eq:datum-adjacent-rows}.  Restriction from \((0,T_*)\) to \((0,T)\)
gives
\[
 \mathfrak R_{N,M}(u;T)\le C_{N,M}
 \qquad(0<T<T_*),
\]
and hence \eqref{f:eq:datum-rectangle-bound}.
Every term in \(\mathfrak R_{N,M}(u;T)\) is the integral of a nonnegative
function, so monotone exhaustion of \((0,T_*)\) gives
\eqref{f:eq:datum-terminal-monotone-realization}.

Conversely, assume \eqref{f:eq:datum-rectangle-bound}.  The same monotone
exhaustion gives the finite full-interval sum in
\eqref{f:eq:datum-terminal-monotone-realization}.  Taking \(N=n\) and \(M=m\)
yields \eqref{f:eq:datum-adjacent-rows} for each finite \((n,m)\).
Proposition~\ref{f:prop:formal-mixed-recurrences} identifies every entry with
the corresponding distributional derivative of the fixed pressure-normalized
history.
\end{proof}

\begin{corollary}[Critical-height exhaustion and Sobolev recovery]
\label{f:cor:datum-critical-height-completion}
For the datum-generated compatible finite rectangles of
Theorem~\ref{f:thm:datum-finite-blocks}, every completed critical coordinate
belongs to \(L^1(0,T_*)\), and
\begin{equation}
 \mathscr H_k[u_0,\nu,f,T_*]
 :=\norm{R_k}_{L^1(0,T_*)}
 =\frac12\norm{\Lambda^{k+1/2}u}_{L^2(0,T_*;L^2)}^2
 \le\frac12\mathfrak R_{0,k}^*[u_0,\nu,f,T_*]
 \qquad(k\in\Nzero).
 \label{f:eq:datum-completed-height-bound}
\end{equation}
For every \(m\in\Nzero\), the completed heights recover the spatial Sobolev
row through
\begin{equation}
 \norm u_{L^2(0,T_*;H^m)}^2
 \le
 2^m\left(
  T_*K_f(T_*)^2+2\norm{R_m}_{L^1(0,T_*)}
 \right)
 \le
 2^m\left(
  T_*K_f(T_*)^2+\mathfrak R_{0,m}^*[u_0,\nu,f,T_*]
 \right).
 \label{f:eq:datum-height-to-spatial-column}
\end{equation}
Thus the compatible critical-height exhaustion gives
\begin{equation}
 (R_k)_{k\in\Nzero}\in\prod_{k\in\Nzero}L^1(0,T_*),
 \qquad
 u\in\bigcap_{m\in\Nzero}L^2(0,T_*;H^m(\R^3)).
 \label{f:eq:datum-complete-height-and-spatial-columns}
\end{equation}
\end{corollary}

\begin{samepage}
\begin{proof}
Take the \(n=0,m=k\) finite rectangle in
\eqref{f:eq:independent-adjacent-finite-rows}; equivalently, take \(N=0\) and
\(M=k\) in \eqref{f:eq:datum-terminal-monotone-realization}.  Since
\(\abs\xi^{2k+1}\le\langle\xi\rangle^{2k+2}\),
\begin{align*}
 2\norm{R_k}_{L^1(0,T_*)}
 &=\norm{\Lambda^{k+1/2}u}_{L^2(0,T_*;L^2)}^2\\
 &\le\norm u_{L^2(0,T_*;H^{k+1})}^2
 \le\mathfrak R_{0,k}^*[u_0,\nu,f,T_*],
\end{align*}
which proves \eqref{f:eq:datum-completed-height-bound}.  For
\(\abs\xi\le1\),
\(\langle\xi\rangle^{2m}\le2^m\); for \(\abs\xi>1\),
\(\langle\xi\rangle^{2m}\le2^m\abs\xi^{2m+1}\).  Hence
\[
 \norm v_{H^m}^2
 \le2^m\left(\norm v_2^2+
 \norm{\Lambda^{m+1/2}v}_2^2\right).
\]
Integrating this inequality, using
\eqref{f:eq:energy-Linf} and
\eqref{f:eq:datum-completed-height-bound}, proves
\eqref{f:eq:datum-height-to-spatial-column}.  Since \(k\) and \(m\) are
arbitrary finite indices, \eqref{f:eq:datum-complete-height-and-spatial-columns}
follows.
\end{proof}
\end{samepage}

\begin{proposition}[Forced spatial-column reconstruction of the mixed intersection]
\label{f:prop:forced-spatial-column-reconstruction}
Let \((T_*,u,p)\) be the pressure-normalized forced history generated by
\((u_0,\nu,f)\), with \(T_*<\infty\), and put \(I=(0,T_*)\).  If
\begin{equation}
 u\in\bigcap_{m\in\Nzero}L^2(I;H^m_\sigma(\R^3)),
 \label{f:eq:forced-complete-spatial-row}
\end{equation}
then \(V_n,Q_n\in C([0,T_*];H^m)\) for every finite \(n,m\), and
\eqref{f:eq:independent-mixed-intersection} follows for this same forced
history.
\end{proposition}

\begin{proof}
Set \(S_m:=\norm u_{L^2(I;H^m)}\).  Sobolev multiplication and the forced
velocity equation give
\begin{equation}
 \norm{u_t}_{L^1(I;H^m)}
 \le \nu\sqrt{T_*}\,S_{m+2}+C_mS_{m+2}^2
     +\norm{\Pp f}_{L^1(I;H^m)}<\infty.
 \label{f:eq:forced-spatial-to-first-temporal-row}
\end{equation}
Thus \(u\in W^{1,1}(I;H^m)\hookrightarrow C([0,T_*];H^m)\), with
\begin{equation}
 \sup_{0\le t\le T_*}\norm{u(t)}_{H^m}
 \le M_{0,m}:=\norm{u_0}_{H^m}
 +\nu\sqrt{T_*}\,S_{m+2}+C_mS_{m+2}^2
 +\norm{\Pp f}_{L^1(I;H^m)}.
 \label{f:eq:forced-spatial-column-base-trace}
\end{equation}

Suppose \(V_0,\ldots,V_n\) are continuous on \([0,T_*]\) in every finite
Sobolev order.  The complete velocity row gives the finite recursive bound
\begin{equation}
 M_{n+1,m}:=\nu M_{n,m+2}
 +C_m\sum_{a=0}^n\binom na M_{a,m+2}M_{n-a,m+2}
 +\sup_{[0,T_*]}\norm{\Pp F_n}_{H^m}<\infty.
 \label{f:eq:forced-spatial-column-temporal-recursion}
\end{equation}
On the interior this continuous expression is \(\partial_tV_n\); integration
identifies its one-sided endpoint derivatives.  Induction gives
\begin{equation}
 V_n\in C([0,T_*];H^m),
 \qquad \sup_{[0,T_*]}\norm{V_n}_{H^m}\le M_{n,m}
 \quad(n,m<\infty).
 \label{f:eq:forced-spatial-column-continuous-rows}
\end{equation}
Simultaneously, the normalized pressure formula
\[
 Q_n=R_iR_j\sum_{a=0}^n\binom na(V_a)_i(V_{n-a})_j
 -(-\Delta)^{-1}\nabla\!\cdot F_n
\]
and \eqref{f:eq:force-pressure-potential-bound} give
\(Q_n\in C([0,T_*];H^m)\) at every finite order.  Since the rows are the
actual temporal derivatives of \(u\),
\begin{equation}
 \norm u_{H^r(I;H^m)}^2
 \le T_*\sum_{n=0}^{r}M_{n,m}^2<\infty.
 \label{f:eq:forced-spatial-column-to-mixed-intersection}
\end{equation}
Taking all finite \((r,m)\) proves
\eqref{f:eq:independent-mixed-intersection}.
\end{proof}

The preceding calculation recovers the same all-orders membership through the
completed heights.  For each finite \(k\), choose a sufficiently high
coordinate rectangle and derive the corresponding critical-height
relationship rectangle on its rows.  This computes
\(R_k=E_{k+1/2}\), with every differentiated force-work term retained in its
recurrence.  Compatibility under enlargement identifies the family
\((R_k)_{k\ge0}\) on the original forced history.  Taking every finite \(k\)
gives \eqref{f:eq:datum-complete-height-and-spatial-columns}.
Proposition~\ref{f:prop:forced-spatial-column-reconstruction} then reconstructs
every temporal and pressure row and recovers
\eqref{f:eq:independent-mixed-intersection}.  Conversely, the
mixed-index exhaustion contains the velocity rows from which every finite
completed height is computed.  Both descriptions use
the same datum-generated block \(\cR_{N,M}[u_0,\nu,f;T_*]\) at each finite
pair and identify the same velocity, pressure, and force derivatives.

\begin{proposition}[Pressure completion]
\label{f:prop:rectangle-pressure-completion}
Under the hypotheses of Theorem~\ref{f:thm:datum-rectangles}, for every
\(N,M\in\Nzero\),
\begin{equation}
 \mathfrak Q_{N,M}^*[u_0,\nu,f,T_*]
 :=\sup_{0<T<T_*}\mathfrak Q_{N,M}(p;T)
 \le
 C_M^{\rm p}(2^{N+1}-1)\mathfrak R_{N,M}^*[u_0,\nu,f,T_*]
 +\sum_{n=0}^N\mathfrak F^{\rm p}_{n,0,M}(0,T_*).
 \label{f:eq:whole-terminal-pressure-bound}
\end{equation}
\begin{samepage}
For every \(0\le n\le N\), the identities
\begin{align}
 Q_n
 &=R_iR_j\sum_{a=0}^n\binom naV_{a,i}V_{n-a,j}
 -(-\Delta)^{-1}\nabla\cdot F_n,
 \label{f:eq:rectangle-pressure-identity}\\
 V_{n+1}
 &=\nu\Delta V_n
 -\sum_{a=0}^n\binom na(V_a\cdot\nabla)V_{n-a}-\nabla Q_n+F_n
 \label{f:eq:rectangle-velocity-identity}
\end{align}
hold respectively in
\(L^1(0,T_*;H^M)\) and \(L^1(0,T_*;H^{M-1})\), including when
\(M=0\).
\end{samepage}
\end{proposition}

\begin{proof}
For every \(0<T<T_*\),
Theorem~\ref{f:thm:formal-finite-row-pressure} gives
\[
 \norm{Q_n}_{L^1(0,T;H^M)}
 +\norm{\nabla Q_n}_{L^1(0,T;H^{M-1})}
 \le C_M^{\rm p}2^n\mathfrak R_{N,M}^*[u_0,\nu,f,T_*]
 +\mathfrak F^{\rm p}_{n,0,M}(0,T).
\]
Sum over \(0\le n\le N\), then let \(T\uparrow T_*\).  Monotone
convergence in the nonnegative norm integrals proves
\eqref{f:eq:whole-terminal-pressure-bound}.  Since
\(\nabla\cdot V_a=0\),
\[
 (V_a\cdot\nabla)V_{n-a}
 =\nabla\cdot(V_{n-a}\otimes V_a).
\]
Lemma~\ref{f:lem:sobolev-product} and Cauchy--Schwarz in time give
\[
 \norm{V_{n-a}\otimes V_a}_{L^1(0,T_*;H^M)}
 \le C_M^{\rm alg}
 \norm{V_a}_{L^2(0,T_*;H^{M+1})}
 \norm{V_{n-a}}_{L^2(0,T_*;H^{M+1})}.
\]
Thus every transport term belongs to
\(L^1(0,T_*;H^{M-1})\), also for \(M=0\).  Moreover,
\[
 \norm{\Delta V_n}_{L^1H^{M-1}}
 \le T_*^{1/2}\norm{V_n}_{L^2H^{M+1}}.
\]
Theorem~\ref{f:thm:datum-rectangles} places \(V_{n+1}\) in
\(L^2(0,T_*;H^{M-1})\), hence in \(L^1(0,T_*;H^{M-1})\).
The prescribed row \(F_n\) belongs to \(L^1(0,T_*;H^{M-1})\), and its
pressure potential belongs to the stated spaces by
\eqref{f:eq:force-pressure-potential-bound}.
The strict-interval identities are identities of distributions; their
full-interval representatives belong to the stated spaces.  This proves the
two identities on \((0,T_*)\).
\end{proof}

\begin{definition}[Finite whole-terminal blocks]
\label{f:def:finite-window-space}
For \(N,M\in\Nzero\), the full-interval block already defined in
\eqref{f:eq:forced-datum-rectangle-record} has, for \(0\le n\le N\),
\begin{equation}
 \begin{aligned}
 V_n&\in L^2(0,T_*;H^{M+1}),&
 V_{n+1}&\in L^2(0,T_*;H^{M-1}),\\
 Q_n&\in L^1(0,T_*;H^M),&
 \nabla Q_n&\in L^1(0,T_*;H^{M-1}),\\
 F_n&\in C([0,T_*];H^{M-1})
       \subset L^1(0,T_*;H^{M-1}).
 \end{aligned}
 \label{f:eq:whole-terminal-rectangle-element}
\end{equation}
Its finite velocity norm is given by
Theorem~\ref{f:thm:datum-rectangles}, its pressure rows by
Proposition~\ref{f:prop:rectangle-pressure-completion}, and its force rows by
the prescribed force class.  Every row is an actual derivative of the one
forced history generated by \((u_0,\nu,f)\).
\end{definition}

\begin{proposition}[Restriction compatibility]
\label{f:prop:rectangle-compatibility}
For \(N'\le N\) and \(M'\le M\),
\begin{equation}
 \left.\cR_{N,M}[u_0,\nu,f;T_*]\right|_{(N',M')}
 =\cR_{N',M'}[u_0,\nu,f;T_*].
 \label{f:eq:rectangle-compatibility}
\end{equation}
Here restriction deletes the temporal rows with index larger than \(N'\) and
reads each retained distribution in the lower Sobolev order \(M'\).  If
\((N'',M'')\le(N',M')\le(N,M)\), two successive restrictions retain exactly
the same coordinates as restriction directly to \((N'',M'')\).  Thus the
datum-generated blocks are a compatible family of literal coordinates of the
same forced history.
\end{proposition}

\begin{proof}
All retained entries are the fixed distributions
\(\partial_t^nu\), \(\partial_t^{n+1}u\), \(\partial_t^np\), and
\(\partial_t^nf\).
The Sobolev embeddings preserve those distributions.  Deleting rows in one
or two stages leaves the same tuple.
\end{proof}
